\documentclass[11pt,a4paper]{article}
\usepackage[utf8]{inputenc}
\usepackage[T1]{fontenc}
\usepackage[spanish,provide=*]{babel}
\addto\captionsspanish{\renewcommand{\tablename}{Tabla}}
\usepackage{lmodern}
\usepackage{microtype}
\usepackage[margin=2.5cm]{geometry}
\usepackage{amsmath,amssymb}
\usepackage{graphicx}
\usepackage{booktabs}
\usepackage{adjustbox}
\usepackage{pdflscape}
\usepackage{float}
\usepackage[font=small,labelfont=bf]{caption}
\usepackage{enumitem}
\usepackage{xcolor}
\usepackage{listings}
\usepackage[round]{natbib}
\usepackage[hidelinks]{hyperref}
\setlength{\parskip}{0.5em}
\setlength{\parindent}{0pt}
\setlist{itemsep=0.2em}

% Texto en rojo = algo que hay que revisar/actualizar después de volver a correr el do-file.
% Para sacar el color: \renewcommand{\revisar}[1]{#1}
\newcommand{\revisar}[1]{\textcolor{red}{#1}}
% Inserta una tabla exportada desde Stata (se achica si no entra)
\newcommand{\tabla}[1]{\begin{adjustbox}{max width=\linewidth, max totalheight=0.72\textheight}\input{#1}\end{adjustbox}}
% Figura que todavía no existe (se genera al correr el do-file): muestra un recuadro
\newcommand{\figura}[2][0.75\linewidth]{\IfFileExists{#2}{\includegraphics[width=#1]{#2}}{\fbox{\parbox{\dimexpr#1-2\fboxsep-2\fboxrule\relax}{\centering\vspace{2em}\textcolor{red}{Figura pendiente: se genera al correr el do-file}\\[0.5em]\texttt{\detokenize{#2}}\vspace{2em}}}}}
% Notas al pie de tablas
\newcommand{\notas}[1]{\par\vspace{0.3em}\begin{minipage}{0.92\linewidth}\footnotesize\textit{Notas:} #1\end{minipage}}
% Consignas en cursiva
\newcommand{\consigna}[1]{\par\textit{#1}\par}

\lstset{
  language={},
  basicstyle=\ttfamily\scriptsize,
  breaklines=true,
  columns=fullflexible,
  keepspaces=true,
  showstringspaces=false,
  frame=single,
  rulecolor=\color{gray!50},
  commentstyle=\color{gray},
  morecomment=[l]{*\ },
  morecomment=[l]{//},
  morecomment=[s]{/*}{*/},
  tabsize=4,
  inputencoding=utf8,
  extendedchars=true,
  literate={á}{{\'a}}1 {é}{{\'e}}1 {í}{{\'i}}1 {ó}{{\'o}}1 {ú}{{\'u}}1
           {Á}{{\'A}}1 {É}{{\'E}}1 {Í}{{\'I}}1 {Ó}{{\'O}}1 {Ú}{{\'U}}1
           {ñ}{{\~n}}1 {Ñ}{{\~N}}1 {ü}{{\"u}}1 {Ü}{{\"U}}1
           {¿}{{?`}}1 {¡}{{!`}}1 {’}{{'}}1 {‘}{{`}}1 {“}{{``}}1 {”}{{''}}1
           {ã}{{\~a}}1 {õ}{{\~o}}1 {â}{{\^a}}1 {ê}{{\^e}}1 {ô}{{\^o}}1 {à}{{\`a}}1
           {Ã}{{\~A}}1 {Õ}{{\~O}}1 {Â}{{\^A}}1 {Ê}{{\^E}}1 {Ô}{{\^O}}1 {Ç}{{\c C}}1 {×}{{$\times$}}1
           {°}{{$^\circ$}}1 {–}{{--}}1 {—}{{---}}1 {·}{{$\cdot$}}1 {ç}{{\c c}}1
}
\begin{document}

\begin{center}
{\Large\bfseries Problem Set 3: Sesgos e imprecisión}\\[0.5em]
E178: Economía Aplicada --- Universidad de San Andrés\\
Fecha de entrega: 6 de septiembre de 2024\\[0.5em]
\textbf{Integrantes:} Juan Batalla, Nina Di Costanzo Pereira, Catalina Gallardo e Ignacio Yenarópulos
\end{center}

\section*{Ejercicio 1}

Trabajamos con un modelo propio. El proceso generador de datos del modelo base es
\[
\text{Wage}_i = 3\,\text{Intelligence}_i + 1\,A_i + 2\,B_i + u_i,
\]
con $n=100$, $\text{Intelligence}\sim N(100,20^2)$, $A\sim N(10,3^2)$, $B\sim N(5,1)$ y $u\sim N(0,1)$ (todas redondeadas a enteros con \texttt{int()}). En cada punto modificamos un único elemento del modelo y comparamos la regresión resultante con la del modelo base. Como los verdaderos coeficientes son conocidos (3, 1 y 2), podemos evaluar tanto el sesgo como la precisión de las estimaciones. Recordemos que, bajo homocedasticidad, la varianza del estimador de MCO de cada coeficiente es
\[
V(\hat\beta_j\mid X)=\frac{\sigma^2}{SST_j\,(1-R^2_j)},
\]
donde $\sigma^2$ es la varianza del error, $SST_j=\sum_i(x_{ij}-\bar x_j)^2$ y $R^2_j$ es el $R^2$ de regresar $x_j$ contra los demás regresores. Esta fórmula explica los resultados de los puntos 1, 2, 3 y 6.

\subsection*{1. ¿Qué sucede con los errores estándar si aumenta el tamaño muestral?}

Disminuyen. Al aumentar $n$, $SST_j$ crece aproximadamente en proporción a $n$, de modo que el error estándar cae a una tasa de $1/\sqrt{n}$. En la Tabla~\ref{tab:1}, al pasar de 100 a 1000 observaciones los errores estándar se reducen aproximadamente a un tercio ($1/\sqrt{10}\approx0.32$): el de \textit{Intelligence} pasa de 0.0036 a 0.0010, el de $A$ de 0.025 a 0.007 y el de $B$ de 0.060 a 0.019. Los coeficientes estimados, a su vez, quedan aún más cerca de sus valores verdaderos.

\begin{table}[H]\centering
\caption{Aumento del tamaño muestral}\label{tab:1}
\tabla{"../output/Tabla 1.tex"}
\notas{Columna (1): modelo base ($n=100$). Columna (2): $n=1000$. Errores estándar entre paréntesis. * $p<0.05$, ** $p<0.01$, *** $p<0.001$.}
\end{table}

\subsection*{2. ¿Qué sucede con los errores estándar si aumenta la varianza de $\mu$?}

Aumentan, porque $\sigma^2$ está en el numerador de la varianza del estimador: con más ruido en el término de error, los datos contienen menos información sobre los parámetros. En la Tabla~\ref{tab:2}, al pasar el desvío del error de 1 a 10, el error estándar de \textit{Intelligence} pasa de 0.0036 a 0.051, el de $A$ de 0.025 a 0.351 y el de $B$ de 0.060 a 0.848. Las estimaciones puntuales también se alejan más de los valores verdaderos (por ejemplo, 0.855 para $A$ y 2.849 para $B$), aunque los estimadores siguen siendo insesgados. El coeficiente de $A$ pasa a ser significativo solo al 5\%.

\begin{table}[H]\centering
\caption{Aumento de la varianza del término de error}\label{tab:2}
\tabla{"../output/Tabla 2.tex"}
\notas{Columna (1): modelo base, $u\sim N(0,1)$. Columna (2): $u\sim N(0,10^2)$. Errores estándar entre paréntesis. * $p<0.05$, ** $p<0.01$, *** $p<0.001$.}
\end{table}

\subsection*{3. ¿Qué sucede con el error estándar de un regresor si aumenta la varianza de $X$?}

Disminuye, porque $SST_j$ está en el denominador: cuanto más varía el regresor, más información hay para identificar su efecto sobre $Y$. En la Tabla~\ref{tab:3} aumentamos el desvío de \textit{Intelligence} de 20 a 50 y su error estándar cae de 0.0036 a 0.0015, en la proporción $20/50=0.4$. Los errores estándar de $A$ y $B$ prácticamente no cambian, porque su variabilidad no se modificó.

\begin{table}[H]\centering
\caption{Aumento de la varianza de un regresor}\label{tab:3}
\tabla{"../output/Tabla 3.tex"}
\notas{Columna (1): modelo base, $\text{Intelligence}\sim N(100,20^2)$. Columna (2): $\text{Intelligence}\sim N(100,50^2)$. Errores estándar entre paréntesis. * $p<0.05$, ** $p<0.01$, *** $p<0.001$.}
\end{table}

\subsection*{4. ¿Cuánto vale la suma de los residuos?}

Si el modelo incluye una constante, la suma de los residuos de MCO es exactamente cero. Es una de las condiciones de primer orden del problema de minimización: la derivada de $\sum_i\hat u_i^2$ respecto de la constante es $-2\sum_i\hat u_i=0$. En el modelo base la suma es $1.04\times10^{-7}$ (Tabla~\ref{tab:4}), que es cero salvo por errores de redondeo numérico.

\begin{table}[H]\centering
\caption{Suma de los residuos del modelo base}\label{tab:4}
\tabla{"../output/Tabla 4.tex"}
\end{table}

\subsection*{5. ¿Son los residuos ortogonales a los regresores?}

Sí. Las condiciones de primer orden de MCO (las ecuaciones normales) son $X'\hat u=0$, es decir, $\sum_i x_{ij}\hat u_i=0$ para cada regresor $j$. Como además los residuos tienen media cero, la correlación muestral entre los residuos y cada regresor es exactamente cero. La Tabla~\ref{tab:5} lo confirma: las correlaciones de los residuos con \textit{Intelligence}, $A$ y $B$ son 0.000. Entre los regresores hay correlaciones bajas que surgen solo por azar (por ejemplo, 0.151 entre $A$ y $B$), ya que los generamos de forma independiente.

\begin{table}[H]\centering
\caption{Correlaciones entre los residuos y los regresores}\label{tab:5}
\tabla{"../output/Tabla 5.tex"}
\end{table}

\subsection*{6. ¿Cómo afecta la alta multicolinealidad a la estimación de $Y$?}

Para estudiarlo generamos \textit{Education} = \textit{Intelligence} + $N(0,1)$, que tiene una correlación de casi 1 con \textit{Intelligence} pero no entra en el verdadero modelo de \textit{Wage}, y la agregamos a la regresión (Tabla~\ref{tab:6}). La multicolinealidad no genera sesgo, pero aumenta mucho la varianza de los coeficientes de las variables colineales, porque $R^2_j$ se acerca a 1: el error estándar de \textit{Intelligence} pasa de 0.0036 a 0.074, unas 20 veces más grande. Lo que el modelo no puede hacer es separar el efecto de cada una de las dos variables colineales.

En cambio, la estimación (predicción) de $Y$ prácticamente no se ve afectada. Como \textit{Intelligence} y \textit{Education} se mueven casi juntas, lo que importa para predecir es la suma de sus coeficientes, $2.932+0.072=3.004$, que es casi idéntica al verdadero efecto de 3 y está estimada con mucha precisión. Es decir, la multicolinealidad afecta la precisión de los coeficientes individuales, pero no la capacidad del modelo de predecir $Y$ (siempre que la relación entre las variables colineales se mantenga).

\begin{table}[H]\centering
\caption{Alta multicolinealidad}\label{tab:6}
\tabla{"../output/Tabla 6.tex"}
\notas{Columna (1): modelo base. Columna (2): modelo que agrega \textit{Education}, con correlación cercana a 1 con \textit{Intelligence}. Como generar \textit{Education} usa números aleatorios adicionales, la muestra de la columna (2) no es idéntica a la del modelo base. Errores estándar entre paréntesis. * $p<0.05$, ** $p<0.01$, *** $p<0.001$.}
\end{table}

\subsection*{7. Error no aleatorio y error aleatorio en $X$}

\textit{Error no aleatorio.} Medimos la inteligencia con un error sistemático constante: $\widetilde{\text{Intelligence}}=\text{Intelligence}+5$. Como $\text{Wage}=3(\widetilde{\text{Intelligence}}-5)+A+2B+u=-15+3\,\widetilde{\text{Intelligence}}+A+2B+u$, el error constante solo se traslada a la constante: los coeficientes de las pendientes y sus errores estándar son idénticos a los del modelo base, y la constante cae exactamente en $3\times5=15$ (de $-0.06$ a $-15.07$; columna 2 de la Tabla~\ref{tab:7}). Si el error sistemático no fuera constante (por ejemplo, si estuviera correlacionado con otras variables), sí sesgaría las pendientes.

\textit{Error aleatorio.} Si $\widetilde{\text{Intelligence}}=\text{Intelligence}+e$ con $e\sim N(0,5^2)$ independiente de todo lo demás (error de medición clásico), el estimador de MCO es sesgado e inconsistente. El regresor observado queda correlacionado con el término de error compuesto ($u-3e$) y el coeficiente se sesga hacia cero (sesgo de atenuación):
\[
\text{plim}\,\hat\beta_1=\beta_1\,\frac{\sigma^2_{\text{Int}}}{\sigma^2_{\text{Int}}+\sigma^2_e}=3\times\frac{400}{425}\approx2.82.
\]
En la columna 3 el coeficiente estimado es 2.741, menor que 3. Además, como el error de medición pasa al término de error, la varianza residual aumenta mucho: los errores estándar crecen, $A$ y $B$ dejan de ser significativos y la constante también queda sesgada (34.92).

\begin{table}[H]\centering
\caption{Error no aleatorio y error aleatorio en el regresor}\label{tab:7}
\tabla{"../output/Tabla 7.tex"}
\notas{Columna (1): modelo base. Columna (2): \textit{Intelligence} medida con error constante (+5). Columna (3): \textit{Intelligence} medida con error aleatorio $N(0,5^2)$. Errores estándar entre paréntesis. * $p<0.05$, ** $p<0.01$, *** $p<0.001$.}
\end{table}

\subsection*{8. Error no aleatorio y error aleatorio en $Y$}

\textit{Error no aleatorio.} Si la variable dependiente se mide con un error constante ($\widetilde{\text{Wage}}=\text{Wage}+10$), el error vuelve a absorberse en la constante. Las pendientes y sus errores estándar no cambian y la constante aumenta exactamente en 10 (de $-0.06$ a 9.94; columna 2 de la Tabla~\ref{tab:8}).

\textit{Error aleatorio.} Si $\widetilde{\text{Wage}}=\text{Wage}+v$ con $v\sim N(0,5^2)$ independiente de los regresores, el error de medición simplemente se suma al término de error. Los estimadores siguen siendo insesgados y consistentes, pero menos precisos, porque aumenta la varianza del error. En la columna 3 los coeficientes siguen cerca de los verdaderos valores (3.027, 1.088 y 2.544, este último dentro de su intervalo de confianza), pero los errores estándar son mucho más grandes (por ejemplo, 0.027 contra 0.0036 para \textit{Intelligence}). A diferencia del error en $X$, el error de medición aleatorio en $Y$ no genera sesgo.

\begin{table}[H]\centering
\caption{Error no aleatorio y error aleatorio en la variable dependiente}\label{tab:8}
\tabla{"../output/Tabla 8.tex"}
\notas{Columna (1): modelo base. Columna (2): \textit{Wage} medido con error constante (+10). Columna (3): \textit{Wage} medido con error aleatorio $N(0,5^2)$. Errores estándar entre paréntesis. * $p<0.05$, ** $p<0.01$, *** $p<0.001$.}
\end{table}

\section*{Ejercicio 2}

Comparamos $\tilde\beta_1$, de la regresión simple de \textit{score} en \textit{attend}, con $\hat\beta_1$, de la regresión que además controla por \textit{cgpa} y \textit{study}. La relación entre ambos está dada por la fórmula del sesgo por variable omitida:
\[
\tilde\beta_1=\hat\beta_1+\hat\beta_2\,\tilde\delta_2+\hat\beta_3\,\tilde\delta_3,
\]
donde $\tilde\delta_j$ es el coeficiente de regresar $X_j$ contra $X_1$. Es decir, $\tilde\beta_1$ y $\hat\beta_1$ difieren solo si las variables omitidas afectan a $Y$ ($\hat\beta_j\neq0$) \emph{y} están correlacionadas con $X_1$ ($\tilde\delta_j\neq0$).

\subsection*{1. $X_1$ muy correlacionada con $X_2$ y $X_3$, que tienen mucha relación con $Y$}

Esperamos que sean distintos. Se cumplen las dos condiciones para que haya sesgo por variable omitida: $\tilde\beta_1$ capta, además del efecto de la asistencia, parte del efecto del promedio y de las horas de estudio. Si, como es razonable, la asistencia está positivamente correlacionada con el promedio y con las horas de estudio, y estas tienen un efecto positivo sobre la nota, $\tilde\beta_1$ sobreestima el efecto de la asistencia ($\tilde\beta_1>\hat\beta_1$).

\subsection*{2. $X_1$ no correlacionada con $X_2$ y $X_3$, pero $X_2$ y $X_3$ muy correlacionadas entre sí}

Esperamos que sean similares. Como $X_1$ no está correlacionada con $X_2$ ni con $X_3$, $\tilde\delta_2\approx\tilde\delta_3\approx0$ y omitirlas no sesga el coeficiente de la asistencia, aunque afecten a $Y$. La alta correlación entre $X_2$ y $X_3$ solo vuelve imprecisos los coeficientes de esas dos variables, pero no afecta a $\hat\beta_1$.

\subsection*{3. $\dot\beta_1$: se agrega el consumo de chocolate ($X_4$)}

Es razonable suponer que, una vez que controlamos por asistencia, promedio y horas de estudio, el consumo de chocolate no afecta la nota de matemática. En ese caso $X_4$ es una variable irrelevante y su inclusión no sesga los demás coeficientes, por lo que esperamos $\dot\beta_1\approx\hat\beta_1$. Esto es así aunque el chocolate esté correlacionado con la asistencia: para que $\hat\beta_1$ estuviera sesgado por omitir $X_4$, el chocolate tendría que afectar a $Y$ \emph{y} estar correlacionado con $X_1$. Solo en ese caso (poco plausible) esperaríamos $\dot\beta_1\neq\hat\beta_1$.

\subsection*{4. $X_1$ muy correlacionada con $X_2$ y $X_3$, que tienen poca relación con $Y$}

Esperamos que sean similares. Aunque $\tilde\delta_2$ y $\tilde\delta_3$ son grandes, $\hat\beta_2$ y $\hat\beta_3$ son cercanos a cero, así que el sesgo por omitirlas es pequeño. La diferencia está en la precisión: en la regresión múltiple, $R^2_1$ (el $R^2$ de regresar $X_1$ contra $X_2$ y $X_3$) es alto, lo que aumenta la varianza de $\hat\beta_1$, y como $X_2$ y $X_3$ casi no explican $Y$, tampoco reducen la varianza del error. Por eso $\hat\beta_1$ será bastante menos preciso que $\tilde\beta_1$.

\subsection*{5. Errores estándar: $X_1$ no correlacionada con $X_2$ y $X_3$, que tienen grandes efectos sobre $Y$ y están muy correlacionadas entre sí}

Esperamos que el error estándar de $\hat\beta_1$ sea menor que el de $\tilde\beta_1$. Al incluir $X_2$ y $X_3$, que explican mucho de $Y$, cae la varianza del error ($\sigma^2$, en el numerador). Como $X_1$ no está correlacionada con ellas, $R^2_1\approx0$ y no hay costo por multicolinealidad para $\hat\beta_1$. La alta correlación entre $X_2$ y $X_3$ solo infla los errores estándar de $\hat\beta_2$ y $\hat\beta_3$. Ambos estimadores son insesgados (punto 2), pero $\hat\beta_1$ es más preciso.

\subsection*{6. Errores estándar de $\dot\beta_1$ y $\hat\beta_1$}

Depende de los supuestos sobre el chocolate:
\begin{itemize}
  \item Si el chocolate no afecta a $Y$ y no está correlacionado con la asistencia (el caso más razonable), los errores estándar serán prácticamente iguales; el de $\dot\beta_1$ será apenas mayor por la pérdida de un grado de libertad.
  \item Si el chocolate no afecta a $Y$ pero está correlacionado con la asistencia (o con $X_2$ y $X_3$), incluirlo aumenta $R^2_1$ sin reducir $\sigma^2$, de modo que el error estándar de $\dot\beta_1$ será mayor que el de $\hat\beta_1$.
  \item Si el chocolate sí afecta a $Y$ y no está correlacionado con la asistencia, incluirlo reduce $\sigma^2$ y el error estándar de $\dot\beta_1$ será menor.
  \item Si el chocolate afecta a $Y$ y está correlacionado con la asistencia, hay dos efectos opuestos (menor $\sigma^2$, mayor $R^2_1$) y el resultado es ambiguo; además, en ese caso $\hat\beta_1$ estaría sesgado.
\end{itemize}

\newpage
\appendix
\section*{Apéndice: código}
\lstinputlisting{../programs/ps3.do}

\end{document}
